\documentclass[10pt]{article}
\providecommand{\FIGDIR}{fig}
\newcommand{\DIGESTDATE}{27 August 2026}
% ===================== PM Research Watch - shared preamble =====================
\usepackage[T1]{fontenc}
\usepackage[utf8]{inputenc}
\usepackage{amsmath,amssymb}
\usepackage{tgpagella}
\usepackage{tgheros}
\renewcommand{\sfdefault}{qhv}
\renewcommand{\ttdefault}{lmtt}
\usepackage{microtype}
\usepackage[english]{babel}

\usepackage[a4paper,top=17mm,bottom=17mm,left=17mm,right=17mm,
            headsep=6mm,footskip=10mm]{geometry}
\usepackage{graphicx}
\usepackage{xcolor}
\usepackage{colortbl}
\usepackage{tikz}
\usetikzlibrary{calc}
\usepackage[most]{tcolorbox}
\usepackage{titlesec}
\usepackage{fancyhdr}
\usepackage{booktabs}
\usepackage{array}
\usepackage{longtable}
\usepackage{enumitem}
\usepackage{multicol}
\usepackage{ragged2e}
\usepackage[hidelinks]{hyperref}

% ------------------------------------------------------------------ palette
\definecolor{Ink}    {HTML}{14262E}
\definecolor{Deep}   {HTML}{0F4C5C}
\definecolor{Teal}   {HTML}{1F7A8C}
\definecolor{Sky}    {HTML}{4FA3C4}
\definecolor{Sage}   {HTML}{6FA287}
\definecolor{Amber}  {HTML}{E3A008}
\definecolor{Clay}   {HTML}{B5643C}
\definecolor{Coral}  {HTML}{D1495B}
\definecolor{Violet} {HTML}{6A5B9E}
\definecolor{Slate}  {HTML}{4A6572}
\definecolor{Paper}  {HTML}{FBF9F5}
\definecolor{Mist}   {HTML}{EEEAE2}
\definecolor{Rule}   {HTML}{D9D3C9}

\pagecolor{Paper}
\color{Ink}

\hypersetup{colorlinks=true, linkcolor=Deep, urlcolor=Teal, citecolor=Teal,
            pdftitle={PM Research Watch}, pdfauthor={Automated daily digest}}

% ------------------------------------------------------------------ headings
\titleformat{\section}
  {\sffamily\bfseries\large\color{Deep}}{}{0pt}
  {}[{\vspace{1.2mm}\color{Amber}\titlerule[1.1pt]}]
\titlespacing*{\section}{0pt}{5.5mm}{2.8mm}

\titleformat{\subsection}
  {\sffamily\bfseries\normalsize\color{Teal}}{}{0pt}{}
\titlespacing*{\subsection}{0pt}{3.6mm}{1.4mm}

\setlist[itemize]{leftmargin=4.2mm, itemsep=1.1mm, topsep=1.3mm, parsep=0pt,
                  label=\textcolor{Teal}{\footnotesize$\blacktriangleright$}}

% ------------------------------------------------------------------ header/footer
\pagestyle{fancy}
\fancyhf{}
\renewcommand{\headrulewidth}{0pt}
\renewcommand{\footrulewidth}{0pt}
\fancyhead[L]{\sffamily\scriptsize\color{Slate}PM RESEARCH WATCH\;\textcolor{Rule}{\textbar}\;\DIGESTDATE}
\fancyhead[R]{\sffamily\scriptsize\color{Slate}Particulate matter monitoring \& health effects}
\fancyfoot[C]{\sffamily\scriptsize\color{Slate}\thepage}

% ------------------------------------------------------------------ components
% metric tile
\newcommand{\metric}[3]{%
  \begin{tcolorbox}[enhanced, colback=white, colframe=Rule, boxrule=0.5pt,
    arc=1.6mm, left=2.2mm, right=2.2mm, top=1.8mm, bottom=1.8mm,
    borderline west={1.7mm}{0pt}{#3}, nobeforeafter, valign=center,
    height=17mm, width=\linewidth]
    \RaggedRight
    {\sffamily\bfseries\fontsize{17}{17}\selectfont\color{#3}#1}\\[0.5mm]
    {\sffamily\fontsize{6.5}{7.4}\selectfont\color{Slate}#2}
  \end{tcolorbox}}

% pull-quote / key insight box
\newtcolorbox{keybox}[1][]{enhanced, colback=Mist, colframe=Mist,
  arc=1.6mm, left=3.2mm, right=3.2mm, top=2.4mm, bottom=2.4mm,
  borderline west={1.7mm}{0pt}{Amber}, #1}

% subtopic band
\newcommand{\band}[2]{%
  \vspace{3.2mm}\noindent
  \begin{tcolorbox}[enhanced, colback=#2, colframe=#2, arc=1.2mm,
    left=2.6mm, right=2.6mm, top=1.1mm, bottom=1.1mm, boxrule=0pt,
    width=\linewidth, nobeforeafter]
    {\sffamily\bfseries\small\color{white}#1}
  \end{tcolorbox}\vspace{1.6mm}}

% one paper entry: {short cite}{title}{journal}{doi}{body}
\newcommand{\paperentry}[5]{%
  \par\vspace{1.9mm}\noindent
  {\sffamily\bfseries\footnotesize\color{Ink}#1}\;%
  {\sffamily\fontsize{7}{8}\selectfont\color{Slate}#3}\\[0.4mm]
  {\footnotesize\itshape\color{Deep}#2}\\[0.6mm]
  {\footnotesize #4}\\[0.3mm]
  {\fontsize{6.6}{7.6}\selectfont\color{Slate}\texttt{doi:}\,\href{https://doi.org/#5}{\texttt{\detokenize{#5}}}}
  \par}

% caption style
\newcommand{\figcap}[1]{\par\vspace{1.0mm}
  {\fontsize{7.2}{8.6}\selectfont\color{Slate}\textsf{\textbf{Fig.}}\; #1\par}\vspace{1.5mm}}

\providecommand{\permil}{\textperthousand}

\setlength{\parindent}{0pt}
\setlength{\parskip}{0pt}
\renewcommand{\arraystretch}{1.18}


\begin{document}

% ==================================================================== MASTHEAD
\thispagestyle{empty}
\begin{tikzpicture}[remember picture, overlay]
  \fill[Deep] (current page.north west) rectangle
      ($(current page.north east)+(0,-67mm)$);
  \fill[Amber] ($(current page.north west)+(0,-67mm)$) rectangle
      ($(current page.north east)+(0,-68.4mm)$);
  \foreach \x/\y/\r in {22/12/0.9, 41/26/1.5, 63/9/0.7, 88/31/1.1, 112/17/1.9,
                        138/29/0.8, 159/13/1.3, 178/54/1.0, 196/21/1.6,
                        30/44/0.6, 74/57/1.2, 125/48/0.9, 168/6/0.8, 191/59/0.7,
                        52/18/1.0, 100/60/0.6, 148/20/1.1, 66/49/1.4, 205/45/0.8}
    \fill[Sky, opacity=0.30] ($(current page.north west)+(\x mm,-\y mm)$) circle (\r mm);
\end{tikzpicture}

\vspace*{4mm}
{\sffamily\fontsize{7.6}{9}\selectfont\color{Sky}%
 AUTOMATED DAILY LITERATURE DIGEST\;\textbullet\;PhD RESEARCH WATCH\par}
\vspace{2.2mm}
{\sffamily\bfseries\fontsize{27}{29}\selectfont\color{white}%
 Particulate Matter\par}
{\sffamily\fontsize{27}{29}\selectfont\color{Sky}%
 Monitoring \& Health Effects\par}
\vspace{5.0mm}
{\sffamily\fontsize{8.4}{10}\selectfont\color{white}%
 \mbox{Issue 2026-08-27}\;\;\textcolor{Amber}{\textbar}\;\;\mbox{Entry window 27 Aug}\;\;
 \textcolor{Amber}{\textbar}\;\;\mbox{36 in scope, 53 out}\par}
\vspace{15mm}

% -------------------------------------------------------------- metric strip
\noindent
\begin{minipage}[t]{0.190\linewidth}\metric{36}{IN-WINDOW RECORDS\\(53 SCREENED OUT)}{Deep}\end{minipage}\hfill
\begin{minipage}[t]{0.190\linewidth}\metric{10}{SUBTOPIC CLUSTERS\\ALL OCCUPIED}{Teal}\end{minipage}\hfill
\begin{minipage}[t]{0.190\linewidth}\metric{13}{INTERVAL-BEARING\\ESTIMATES}{Sky}\end{minipage}\hfill
\begin{minipage}[t]{0.190\linewidth}\metric{14}{TIER-A\\RECORDS}{Amber}\end{minipage}\hfill
\begin{minipage}[t]{0.190\linewidth}\metric{4}{ABSTRACT-FREE\\DEPOSITS (11\%)}{Clay}\end{minipage}

\vspace{5mm}

% -------------------------------------------------------------- signal of day
\section*{\sffamily\bfseries\color{Deep}Signal of the day}
\vspace{-1.5mm}{\color{Amber}\rule{\linewidth}{1.1pt}}\vspace{2.5mm}

\begin{keybox}
{\sffamily\bfseries\small\color{Deep}A meta-analysis of 19 studies shows that the
\emph{harmonisation step} --- not the exposure --- carries most of the heterogeneity in
the PM--venous-thromboembolism literature.}\\[1.4mm]
\footnotesize
Pooling 19 observational studies and more than \textbf{4.5 million participants}, the
headline long-term estimates are null and enormously heterogeneous:
\textbf{PM$_{2.5}$ HR 1.42 (95\,\% CI 0.88--2.28), $I^2 = 97.0\,\%$} and
\textbf{PM$_{10}$ HR 1.18 (0.85--1.64), $I^2 = 96.2\,\%$}. Restricting to the subset of
studies that \emph{natively reported} effects per 10\,$\mu$g/m$^3$ --- i.e.\ dropping
every estimate the reviewers had to rescale --- collapses the heterogeneity to
\textbf{$I^2 = 0\,\%$} and yields \textbf{PM$_{2.5}$ HR 1.070 (1.010--1.135)} and
\textbf{PM$_{10}$ HR 1.023 (1.000--1.047)}. Two orders of magnitude of $I^2$ removed by a
data-handling decision is not a subgroup finding; it says the rescaling of
non-native contrasts (per-IQR, per-SD, per-quartile) into a common per-10 unit is
importing variance that the underlying studies do not contain. The direct consequence
for distributed sensing: a network that reports its own exposure contrast in a
non-standard unit is not merely inconvenient for meta-analysis, it is
\emph{actively destabilising} to it. Reporting per-10\,$\mu$g/m$^3$ natively is the
cheapest methodological contribution a monitoring study can make.
\end{keybox}

\vspace{2mm}

% -------------------------------------------------------------- exec brief
\section*{\sffamily\bfseries\color{Deep}What changed today}
\vspace{-1.5mm}{\color{Amber}\rule{\linewidth}{1.1pt}}\vspace{2.5mm}

\footnotesize
\begin{itemize}\itemsep2pt
\item \textbf{The largest single-day issue this watch has built, and only the fourth
to occupy all ten clusters.} 36 records against a median of roughly 14 --- mechanical,
not scientific: a \textit{Toxics} volume release supplied six, and Crossref's by-ISSN leg
43 deposits of which most were rejected. Only the 26 July founding backfill (164 records
over a month) is larger. \textbf{Four deposits arrived abstract-free}, three of them on
the sensing axis --- the refractive-index, fire/extreme-PM, TRAKER and Yangtze records.

\item \textbf{Low-cost PM sensors used as a biological proxy, with a validated
transfer function.} Li et al.\ treat greenhouse PM sensors not as air-quality
instruments but as continuous physical proxies for airborne \textit{Pseudomonas}
aerosol load, triggering targeted qPCR only when an estimated risk threshold is crossed.
Across \textbf{732 plot-days} the surveillance model reached
\textbf{$R^2 = 0.884$ for pathogen aerosol accumulation} and the downstream forecast
\textbf{$R^2 = 0.874$ for daily disease-index change rate}. The framing is the
interesting part: a sensor with no biological specificity is made useful by pairing it
with a specific but expensive assay that it schedules.

\item \textbf{Condensable PM gets a real-time instrument and an age-resolved dose.}
Kang et al.\ characterised CPM from eight vehicle types with a purpose-built real-time
system (PNU-CPM). CPM accounted for \textbf{63.2\,\% of total PM mass for an untreated
level-4 diesel} against \textbf{8.2\,\% for GDI}; organics were dominated by alkanes
($\sim$53.5\,\%) and the water-soluble fraction by SO$_4^{2-}$ and Na$^+$. MPPD
simulation gives higher mass-per-body-weight lung deposition in younger age groups.
CPM is invisible to hot-filter source testing by construction, and semi-volatile
material of this composition will not obey the scattering and hygroscopic assumptions
baked into optical-sensor corrections.

\item \textbf{Speciation beats mass again, on oxidative potential.} In Xi'an,
XGBoost--SHAP attribution over DTT assays makes \textbf{potassium and manganese} --- a
biomass-burning tracer and a redox-active transition metal --- the dominant contributors
to volume-normalised OP, despite water-soluble metals being a minor mass fraction. Dust
dominates winter OP by loading, traffic dominates summer OP by intrinsic toxicity. Same
conclusion as the 26 August component-IRR result, different measurement chain.

\item \textbf{A well-powered null more informative than most positives.} Instrumenting
municipal PM$_{2.5}$ with wind direction across \textbf{$>$3.5 million French
municipality-weeks}, Morvan et al.\ find \emph{no} effect on COVID-19 hospitalisation or
death among people without pre-existing conditions, and a clear one once comorbidity is
conditioned on (\textbf{+1.5 deaths per 100{,}000 per 10\,$\mu$g/m$^3$}, 0.8--2.2;
respiratory 0.9 (0.4--1.4), cardiovascular 0.6 (0.1--1.1)). The 2020--21 literature was
overwhelmingly ecological and positive; this design has an instrument, and it puts the
effect entirely in the susceptible stratum.

\item \textbf{PM$_1$ outperforms PM$_{2.5}$ in two Guangzhou cohorts.} Prenatal PM$_1$
gave the largest per-5\,$\mu$g/m$^3$ odds of a rapid fetal-weight trajectory
(\textbf{OR 1.56, 1.03--2.37}, versus 1.44 for PM$_{2.5}$ and 1.28 for PM$_{10}$), and
the toddler-sleep cohort put the sub-micron fraction on par with PM$_{2.5}$. One city,
one exposure vendor, so independence is partial --- but the ordering is what deposition
predicts, and what mass-only monitoring cannot see.

\end{itemize}

\clearpage

% ==================================================================== ANALYTICS
\section*{\sffamily\bfseries\color{Deep}Corpus at a glance}
\vspace{-1.5mm}{\color{Amber}\rule{\linewidth}{1.1pt}}\vspace{3mm}

\noindent\includegraphics[width=\linewidth]{\FIGDIR/f1_subtopics.png}
\figcap{Records per subtopic cluster. \textbf{All ten clusters are occupied --- the
fourth issue ever to manage it, and the first since 28 July.}
\textit{Exposure assessment \& modelling} 7 and
\textit{Sensing} 6 lead; \textit{Reproductive \& developmental} is the thinnest at 1. The
completeness is a consequence of volume (36 records), not of breadth in the day's
deposits: at 14 records this distribution would have left three or four clusters empty.}

\vspace{2mm}
\noindent\includegraphics[width=\linewidth]{\FIGDIR/f2_design_endpoint.png}
\figcap{Left --- study architecture. \textit{Measurement campaign} 8,
\textit{Observational --- cohort} 6, \textit{Modelling / inventory} 5,
\textit{Review / synthesis} 4, \textit{Metadata only} 4,
\textit{Observational --- cross-sectional} 3, \textit{Ecological} 3,
\textit{Experimental / toxicology} 2, \textit{Observational --- acute} 1.
The \textit{Metadata only} bar is a harvesting artefact, not a design: those four records
have no abstract and their architecture is unknown. Right --- health endpoint.
\textbf{Twelve records carry no human health endpoint} (instrument or exposure work);
\textit{Mortality} 5 and \textit{Respiratory} 4 lead the remainder, and seven endpoint
categories are singletons.}

\vspace{2mm}
\noindent\includegraphics[width=\linewidth]{\FIGDIR/f3_metric_geography.png}
\figcap{Left --- particle metric. \textit{PM$_{2.5}$ only} 13,
\textit{PM (unspeciated) / emissions} 11, \textit{PM$_{2.5}$ + PM$_{10}$ jointly} 8,
\textit{ultrafine / nanoscale} 2, and one each for \textit{composition / speciation} and
\textit{size-resolved / number}. The single \textit{size-resolved} record is the PFAS
worker panel, whose entire finding depends on the size cut. Note the grouper undercounts
speciation: the Xi'an OP study, the New Mexico metals study and the Guangzhou toddler
study all resolve constituents but bin as PM$_{2.5}$-only. Right --- geography.
\textbf{Global / multi-region 10 is inflated} --- it absorbs both genuine multi-region
syntheses and records whose site is not stated, including three of the four abstract-free
deposits. China 8 and Europe 8 dominate the located records; North America contributes 1.}

\vspace{2mm}
\noindent\includegraphics[width=\linewidth]{\FIGDIR/f4_heatmap.png}
\figcap{Subtopic $\times$ study-architecture cross-tabulation. The densest cell is
\textit{Exposure assessment} $\times$ \textit{Measurement campaign} at 4, followed by
\textit{Sensing} $\times$ \textit{Measurement campaign} at 3. \textbf{Five cells hold 2}
--- \textit{Sensing} $\times$ \textit{Metadata only}, \textit{Burden} $\times$
\textit{Modelling / inventory}, \textit{Mechanistic toxicology} $\times$
\textit{Experimental / toxicology}, \textit{Neuro} $\times$
\textit{Observational --- cohort} and \textit{Other clinical endpoints} $\times$
\textit{Observational --- cross-sectional} --- and the remainder are singletons. The diagonal structure is the
expected one --- instrument and exposure work concentrates in measurement campaigns,
clinical clusters in cohorts --- and the only off-diagonal density is the pair of
metadata-only records sitting in \textit{Sensing}, which is an artefact of which
publishers deposit abstracts, not of what was studied.}

\vspace{2mm}
\noindent\includegraphics[width=\linewidth]{\FIGDIR/f5_forest.png}
\figcap{Thirteen interval-bearing estimates from five studies, all on particulate
contrasts. \textbf{Three internally comparable sets:} the three VTE meta-analytic hazard
ratios (same review, but the wide PM$_{2.5}$ 1.42 estimate and the narrow 1.070 estimate
are the \emph{same exposure} under two harmonisation rules and must be read as a pair,
not as two findings); the three per-5\,$\mu$g/m$^3$ fetal-growth odds ratios, which are
mutually comparable and ordered PM$_1 >$ PM$_{2.5} >$ PM$_{10}$; and the five New Mexico
metal rate ratios, which share a cohort, an exposure model and an outcome. Two estimates
are \emph{not} exposure--outcome associations on a per-unit PM contrast: the UK Biobank
RA hazard ratio is per SD of a \emph{metabolic signature} derived from PM$_{2.5}$, not
per unit PM$_{2.5}$, and belongs on this axis only in the weak sense that it shares a
ratio scale. The S\~ao Paulo respiratory RR crosses the null (0.889--1.204) and its
source reports $p = 0.054$; it is drawn because it is a per-10\,$\mu$g/m$^3$ PM$_{2.5}$
contrast, not because it is significant.}

\vspace{2mm}
\noindent\includegraphics[width=\linewidth]{\FIGDIR/f6_lifecourse.png}
\figcap{Life-course windows addressed, counted only where a record's study population is
explicitly anchored to a stage. Six of 36 records qualify. \textbf{All five windows are
occupied --- the first fully populated strip since 19 August}, and only the seventh in
this watch's history. The paediatric window breaks a run of \textbf{four} consecutive
empty issues and the adolescent window a run of \textbf{five}. Working-age adulthood
carries two (the Flemish endothelial cohort and the 20-worker PFAS panel); every other
window carries exactly one, so the strip is coverage, not weight of evidence.}

\clearpage

% ==================================================================== DIGEST
\section*{\sffamily\bfseries\color{Deep}Paper-level digest}
\vspace{-1.5mm}{\color{Amber}\rule{\linewidth}{1.1pt}}\vspace{1mm}

{\fontsize{7.4}{9}\selectfont\color{Slate}Ordered by cluster; the 27 August entry window
only. Tier A/B/C is a triage of how far the record can be relied on, not of journal
prestige.\par}

\band{Sensing, forecasting \& instrumentation\hfill 6 records}{Amber}

\paperentry{Li et al.\ 2026c}
{Early warning of cucumber angular leaf spot from airborne pathogen aerosols estimated with PM sensors}
{Plants (Basel)}
{Two-stage greenhouse surveillance in which \textbf{low-cost PM sensors act as continuous
physical proxies} for suspended pathogen aerosol, and qPCR air sampling is triggered only
when an estimated aerosol risk crosses a threshold. Evaluated on \textbf{732 plot-days}
across spatially independent trials: an Extra Trees model predicted pathogen-aerosol
accumulation at \textbf{$R^2 = 0.884$}, and an XGBoost model integrating aerosol load with
microclimate forecast the daily disease-index change rate at \textbf{$R^2 = 0.874$}. SHAP
identified aerosol concentration and vapour pressure deficit as the primary drivers.
The limitation is stated by the authors and is the right one: PM sensors have no
biological specificity, and the design compensates by scheduling a specific assay rather
than by claiming the sensor sees pathogens. Cross-greenhouse rather than cross-site
validation, and no reference-instrument comparison for the PM channel itself.}
{10.3390/plants15162510}

\paperentry{Kang et al.\ 2026b}
{Emission characteristics and health impacts of vehicle-originated condensable PM}
{J Hazard Mater}
{Condensable PM from \textbf{eight vehicle types} characterised with a purpose-built
real-time system (PNU-CPM) plus offline physicochemical analysis, with emission factors
derived from the real-time record. Untreated \textbf{diesel levels 4 and 5 emitted the
most CPM and total PM}; CPM mass fractions were \textbf{63.2\,\% (diesel L4), 8.2\,\%
(GDI) and 58.8\,\% (LPG 2)}. Organic CPM was $\sim$53.5\,\% alkanes and $\sim$5.1\,\%
esters; dominant water-soluble ions were SO$_4^{2-}$ and Na$^+$. MPPD dosimetry gives
higher deposited mass per body weight in younger age groups. Directly relevant to optical
sensing: CPM is ultrafine, semi-volatile and forms after the tailpipe, so it is missed by
hot-filter source testing and will not follow the scattering or hygroscopic-growth
assumptions that low-cost sensor corrections inherit from combustion soot. No on-road
fleet measurement --- these are dynamometer emission factors.}
{10.1016/j.jhazmat.2026.143367}

\paperentry{Lefevre et al. 2026}
{Urban microclimate dataset from low-cost sensor networks, Reunion Island}
{Data Brief}
{A data descriptor for \textbf{25 self-built LoRa weather stations} deployed on
residential buildings across two contrasting neighbourhoods (Saint-Denis, Le Port),
October 2024 to June 2025, at \textbf{10-minute cadence} for air temperature, relative
humidity, wind speed, global horizontal solar radiation and grey-globe temperature.
Station metadata, QC procedures and maintenance logs are published alongside, and
incomplete stations are identifiable. \textbf{There is no PM channel} --- this is a
microclimate network, not an air-quality one --- so its value here is as a deployment
reference: dense neighbourhood-scale LoRa siting, self-built hardware, published QC, and
an explicitly stated reuse path for model calibration (ENVI-met, STEVE). Tropical
hot-humid deployments are exactly the regime in which PM sensor RH correction fails
worst, and this is a rare open record of what such a network's data actually look like.}
{10.1016/j.dib.2026.113166}

\paperentry{Wang et al.\ 2026d}
{Perchloric acid drives new particle formation enhanced by dimethylamine, ammonia and sulfuric acid}
{Atmos Chem Phys}
{Quantum-chemical calculations coupled to the Atmospheric Clusters Dynamic Code compare
the nucleation enhancement of DMA, NH$_3$ and H$_2$SO$_4$ for perchloric-acid-based new
particle formation, motivated by recent Arctic observations of atmospheric HClO$_4$.
Both DMA and NH$_3$ bind PA strongly through hydrogen bonding and proton transfer, and
\textbf{PA nucleates more readily with DMA than sulfuric acid does}; the PA--DMA cluster
formation rate exceeds PA--NH$_3$ even at overwhelming NH$_3$ excess, with 1:1
PA:DMA clusters controlling growth. Purely computational --- no chamber or field
verification, and the claimed NH$_3$ excess of ``10--100 orders of magnitude'' is a
transcription problem in the abstract rather than a physical statement. Relevant as a
candidate nucleation channel in clean marine and polar boundary layers where sulfuric
acid is scarce, which is where number-concentration instruments see unexplained events.}
{10.5194/acp-26-12171-2026}

\paperentry{Moteki et al. 2026}
{Refractive index of fractal-like aggregates from complex forward-scattering amplitudes}
{J Aerosol Sci}
{\textbf{Metadata-only deposit} --- title and DOI without an abstract. The method named in
the title is the single-particle complex forward-scattering amplitude $S(0^\circ)$
approach, previously used to constrain the black-carbon refractive index at
$\lambda = 0.633\,\mu$m; the new element is \textbf{mass classification of the aggregates}
before the optical measurement, which removes the size-distribution term that dominates
the uncertainty in bulk retrievals. If it delivers, this is the most directly useful
sensing record in the window: every optical PM instrument, from a nephelometer to a
\$20 laser scatter sensor, converts scattered intensity to mass through an assumed
refractive index, and for soot-dominated aerosol that assumption is the largest single
error term. Carried at tier C on the title alone and placed at the top of the abstract
retry list.}
{10.1016/j.jaerosci.2026.106884}

\paperentry{Chen et al. 2026e}
{Aerosol and VOC pollution characteristics along the lower Yangtze, from shipborne observation}
{Atmos Environ}
{\textbf{Metadata-only deposit.} Shipborne mobile measurement of aerosol and VOCs along an
inland waterway corridor --- a platform class this watch sees rarely and one with a
specific methodological problem: the measurement frame moves through the plume it is
sampling, so concentrations are Lagrangian-contaminated in a way that fixed-site data are
not, and self-sampling of the vessel's own exhaust is a standing artefact that requires
wind-relative rejection to control. Whether the authors address either is not
determinable from the title. Carried at tier C; the record is retained rather than
rejected because inland-shipping emissions are a genuinely under-monitored PM source and
the corridor is one of the most heavily trafficked in the world.}
{10.1016/j.atmosenv.2026.122331}

\band{Exposure assessment \& modelling\hfill 7 records}{Deep}

\paperentry{Rashid et al. 2026}
{Ultrafine particles and mortality: scoping review of evidence and exposure-assessment limits}
{Toxics}
{PRISMA-ScR scoping review of PubMed and Scopus: \textbf{21 of 704 screened articles},
2007--2025, split \textbf{14 short-term} (time-series or case-crossover against central-site
particle number counters) and \textbf{7 long-term} (land-use regression, chemical transport
or machine-learning surfaces). \textbf{10 of 14 short-term and 6 of 7 long-term studies
reported positive associations}, strongest for respiratory mortality. The finding that
matters for this watch is the second-order one: \textbf{the strongest associations came
from the studies with higher-resolution or source-specific exposure assessment}, a pattern
consistent with reduced Berkson-type misclassification rather than with true effect
heterogeneity. A scoping review does not pool, so there is no summary estimate and no
formal heterogeneity statistic; the exposure-method gradient is observed, not tested.
Still the clearest statement to date of why UFP regulation stalls on measurement.}
{10.3390/toxics14080664}

\paperentry{Chen et al. 2026f}
{Source-specific oxidative potential of PM$_{2.5}$: water-soluble metals via DTT and interpretable ML}
{Toxics}
{Winter and summer ambient PM$_{2.5}$ in Xi'an assayed for DTT oxidative potential, with
XGBoost--SHAP attribution over composition and receptor-model source contributions.
\textbf{Volume-normalised OP (DTT$_v$) was markedly higher in winter}, and although
water-soluble metals are a small fraction of PM$_{2.5}$ mass, \textbf{potassium and
manganese dominated the OP attribution} --- a biomass-burning tracer and a redox-active
transition metal. Source apportionment separates the two seasons cleanly: dust drives
winter OP through mass loading, traffic drives summer OP through intrinsic per-mass
toxicity. The abstract reports no numeric DTT$_v$ values, confidence intervals or model
performance, which is a real gap for a paper whose central claim is quantitative
attribution; SHAP importance is also not a causal decomposition and correlated
constituents will trade rank. The qualitative conclusion --- OP tracks composition and
source, not mass --- is the second independent instance of it this week.}
{10.3390/toxics14080646}

\paperentry{Li et al. 2026d}
{Surface PM transport under a closed earth-air model}
{Environ Pollut}
{One year of continuous in-situ monitoring in the semi-enclosed basement of the Mogao
Grottoes conservation department, testing whether \textbf{earth-air vertical movement}
driven by atmospheric-pressure fluctuation controls surface PM independently of wind and
human activity. Structural equation modelling gives \textbf{CO$_2$ as a tracer for
EVM-driven PM dynamics}, with path coefficients of \textbf{0.452 (PM$_{2.5}$), 0.470
(PM$_1$) and 0.386 (PM$_{100}$)}; $^{222}$Rn regulates PM positively but far more weakly,
and relative humidity is strongly inhibitory. PM$_{2.5}$ and PM$_1$ are strongly affected
by EVM, PM$_{100}$ only weakly. Single site, single enclosure, and SEM path coefficients
on autocorrelated one-year time series overstate precision badly. The mechanism claim is
interesting for indoor and heritage monitoring; the effect sizes should not be
transported.}
{10.1016/j.envpol.2026.128983}

\paperentry{Mainka \& Mainka 2026}
{Urban cycling versus PM$_{2.5}$ exposure: does the benefit outweigh the risk?}
{Preprints.org}
{Scenario analysis of \textbf{162 real commuting trips by two cyclists} in Gliwice across
heating and non-heating seasons, with personal PM$_{2.5}$ measured by low-cost sensor and
inhaled dose reconstructed from concentration, trip duration and cycling-specific
ventilation. Median trip 33.0 min/day. The median physical-activity relative risk was
\textbf{0.836} (a 16.4\,\% all-cause mortality reduction) against a median PM$_{2.5}$
relative risk of \textbf{1.00182} (0.182\,\% increase), for a \textbf{net RR of 0.840
(IQR 0.815--0.879)}; \textbf{160 of 162 trips retained net benefit} under sensitivity
analysis. Two hard limits: \textbf{$n = 2$ commuters}, so the exposure distribution is a
convenience sample of two routes, and the risk functions are transported from published
cohorts rather than estimated. Unreviewed preprint. The useful part is the
method --- trip-resolved dose from a low-cost sensor is the right exposure object for
this question, and almost no one computes it.}
{10.20944/preprints202608.1778.v1}

\paperentry{Moryani et al. 2026}
{Bioaccessibility and speciation of heavy metals in PM$_{2.5}$ road dust}
{Toxics}
{\textbf{48 grid-based samples} across four functional zones (residential, school,
hospital, office) in Karachi ($>$27 million) and Shikarpur ($\sim$0.3 million), with the
PM$_{2.5}$ fraction of road dust extracted in artificial lysosomal fluid
(\textbf{pH 4.5, 37\,$^\circ$C, 24 h}) and speciation modelled in PHREEQC and Visual
MINTEQ. \textbf{Cr, As and Pb showed moderate inhalation bioaccessibility}; contamination
and ecological-risk indices were computed alongside. The design addresses a real gap ---
total metal content is a poor predictor of the dose that reaches lung tissue --- but ALF
extraction is a simulation with no in-vivo anchor, and the abstract reports no
bioaccessible fractions, no concentrations and no zone contrasts, so nothing quantitative
can be lifted. Road dust is also resuspended, not directly emitted, which the pollution
indices do not distinguish.}
{10.3390/toxics14080704}

\paperentry{Isworo \& Oetari 2026}
{Risk-integrated industrial air quality: Gaussian dispersion linked to a health-weighted pressure index}
{Environ Sci Pollut Res Int}
{Gaussian-plume dispersion for a Central Java textile facility combined with a
hazard-quotient-normalised Atmospheric Environmental Pressure Index, using an emission
inventory, ERA5 2015--2024 meteorology and six receptor measurements. The headline is a
compliance paradox worth stating: \textbf{modelled concentrations meet Indonesia's
national standards while PM$_{2.5}$ exceeds the WHO 2021 guideline}, and PM$_{2.5}$ and
PM$_{10}$ dominate the index. AEPI values (0.06 baseline, 0.08 predicted) sit in the
authors' own ``very low'' band, which makes the index insensitive at exactly the
concentrations where the health argument bites --- a scaling problem the paper does not
confront. One facility, six receptors, no validation of the plume model against the
measurements. Tier C: the framework is transferable, the result is not evidence.}
{10.1007/s11356-026-38159-z}

\paperentry{Wu et al. 2026b}
{Global fire activity exacerbates extreme PM$_{2.5}$ pollution}
{npj Clim Atmos Sci}
{\textbf{Metadata-only deposit.} The title claim --- that fire activity acts
disproportionately on the \emph{extreme tail} of the PM$_{2.5}$ distribution rather than
on the mean --- is the one that matters, because attainment statistics, exposure-response
functions fitted on annual means and most burden calculations are all mean-anchored and
will understate a tail-loading source. It is also the regime in which low-cost networks
are most useful and least characterised: sensor response at several hundred
$\mu$g\,m$^{-3}$ of smoke is outside the concentration range of essentially every
published co-location study. Carried at tier C pending the abstract; a Nature-portfolio
deposit with no abstract in Crossref is unusual and may resolve on retry.}
{10.1038/s41612-026-01527-0}

\band{Cardiovascular \& metabolic\hfill 3 records}{Teal}

\paperentry{Correa Lara et al. 2026}
{Ambient PM and venous thromboembolism: updated systematic review and meta-analysis}
{Thromb Haemost}
{\textbf{19 observational studies, $>$4.5 million participants}, searched across PubMed,
Embase, Web of Science and Cochrane, harmonised to 10\,$\mu$g/m$^3$ and pooled under
restricted maximum likelihood with Hartung--Knapp adjustment. Overall long-term:
\textbf{PM$_{2.5}$ HR 1.42 (0.88--2.28), $I^2 = 97.0\,\%$}; \textbf{PM$_{10}$ HR 1.18
(0.85--1.64), $I^2 = 96.2\,\%$}. Restricting to studies natively reporting per
10\,$\mu$g/m$^3$: \textbf{PM$_{2.5}$ 1.070 (1.010--1.135)} and \textbf{PM$_{10}$ 1.023
(1.000--1.047)}, both at \textbf{$I^2 = 0\,\%$}. Short-term exposures were null. The
sensitivity analysis is the finding --- near-total heterogeneity removed by excluding
rescaled contrasts --- but it is also a smaller, non-random subset, and a subset selected
on a reporting characteristic can be homogeneous for reasons other than correctness.
Inconsistent VTE case definitions across the 19 studies remain unaddressed.}
{10.1055/a-2945-4974}

\paperentry{Yang et al. 2026b}
{Endothelial function at low-level chronic residential air pollution}
{Blood Press}
{\textbf{453 Flemish adults} (47.7\,\% women, mean age 52.8 y) with endothelial function
by finger photoplethysmography after 5 min ischaemia, related to residential PM$_{10}$,
PM$_{2.5}$, NO$_2$ and black carbon from high-resolution spatiotemporal interpolation.
Six-year means were \textbf{15.9, 12.8, 14.3 and 1.04\,$\mu$g/m$^3$} --- at or below
current EU limits and near the WHO interim range. \textbf{Dmax was inversely associated
with PM$_{2.5}$} irrespective of risk-factor adjustment; other pollutant--outcome pairs
were consistently inverse but weaker ($p < 0.10$). Association sizes \textbf{weakened as
pollution fell over 2010--2015} ($p \leq 0.044$), and the spatial gradient in endothelial
dysfunction tracked the PM$_{2.5}$ gradient. Both endothelial measures predicted a
composite cardiovascular endpoint (HR 0.27--0.49). Modest $n$, and no confidence intervals
are given for the exposure--endothelium associations themselves, which is why none appear
in the forest plot. The attenuation-with-decline result is the notable one and is rarely
reported.}
{10.1080/08037051.2026.2724736}

\paperentry{Ram\'irez-Moreno et al. 2026}
{Regional stroke hospitalisation and mortality against long-term PM$_{2.5}$}
{Med Sci (Basel)}
{Nationwide ecological analysis of Spain's 17 autonomous communities, 2013--2021,
\textbf{306 sex-specific region-year observations}, linking INCLASNS stroke indicators to
MITECO ambient pollution. Mean age-standardised hospitalisation rate 20.1 $\pm$ 5.4 per
10{,}000. \textbf{PM$_{2.5}$ had the largest positive correlation with age-standardised
mortality of any pollutant examined ($r = 0.264$, $p < 0.001$)}, which the authors
correctly describe as weak; associations with hospitalisation and with other pollutants
were small and inconsistent in sign. \textbf{Regional variation in stroke burden was not
explained by the available pollution indicators.} Exposure is at region-year resolution
against sex-stratified outcomes, so the sex analysis is uninterpretable by construction
and the authors say so. Reported here as a well-labelled null; the design cannot support
individual-level inference and does not claim to.}
{10.3390/medsci14040475}

\band{Neuro / mental health\hfill 2 records}{Deep}

\paperentry{Wu et al. 2026c}
{Early-life PM exposure and toddler sleep: particle size and constituents}
{Toxics}
{Birth cohort of \textbf{514 mother--child pairs} in Guangzhou followed to age two
(2017--2018), with PM$_1$, PM$_{2.5}$, PM$_{10}$ and PM$_{2.5}$ constituents from
high-resolution ground-level grids and sleep by the Chinese Brief Infant Sleep
Questionnaire. Per-IQR early-life exposure was associated with reduced nighttime sleep
duration: \textbf{PM$_1$ $\beta = -0.18$ (95\,\% CI $-0.32$, $-0.04$), PM$_{2.5}$
$-0.23$ ($-0.38$, $-0.08$), PM$_{10}$ $-0.21$ ($-0.35$, $-0.07$)} hours. Quantile
g-computation on the constituent mixture reproduced the decrement with \textbf{SO$_4^{2-}$
as the dominant contributor}, and pregnancy-window findings were consistent. Outcome is
parent-reported, which biases toward non-differential misclassification but not
necessarily away from confounding by household environment; the three particle sizes are
collinear at this resolution, so the size ordering is suggestive rather than
discriminating. Sulfate as the driver is consistent with the secondary-inorganic signal
this watch carried on 26 August from an unrelated endpoint.}
{10.3390/toxics14080653}

\paperentry{Hendrickx et al. 2026}
{Recent air pollution exposure, children's neurocognition and endothelial function}
{Int J Environ Res Public Health}
{\textbf{138 healthy children} (mean age 10.3 y, 55.8\,\% male) across seven Antwerp
primary schools including one open-air school, with indoor and outdoor PM and outdoor
NO$_2$ measured and endothelial and attentional outcomes assessed at three repeated
visits. Exposure was resolved at \textbf{4 h (direct), 24 h (acute) and 1 week (recent)}
windows. \textbf{Higher direct indoor PM was associated with reduced microvascular
dilatation}, while \textbf{selective and sustained attention \emph{improved} with direct
and recent PM exposure} respectively --- a sign reversal the authors report without
resolving, and which most plausibly reflects the same seasonal or activity pattern
driving both classroom occupancy and outdoor concentration rather than a cognitive
benefit. No NO$_2$ associations. Excluding the open-air school shifted results, which on
$n = 138$ with seven schools means the school-level clustering is doing real work. Tier B
for the vascular finding; the attention result should not be cited without the full text.}
{10.3390/ijerph23081070}

\band{Respiratory \& allergic\hfill 2 records}{Sage}

\paperentry{Morvan et al. 2026}
{Causal effect of PM$_{2.5}$ on COVID-19 hospitalisation and death, conditioned on comorbidity}
{Int J Epidemiol}
{Instrumental-variable analysis over \textbf{more than 3.5 million weekly municipality
observations} covering all of France, 2020--2021, with COVID-19 hospitalisations and
deaths disaggregated by age, sex and municipality and linked to records of pre-existing
conditions. \textbf{Wind direction instruments local PM$_{2.5}$} in a two-stage least
squares design. Among individuals \emph{without} pre-existing conditions there was
\textbf{no significant effect} on either outcome. Conditioning on comorbidity, a
10\,$\mu$g/m$^3$ increase was associated with \textbf{+1.5 COVID-19 deaths per 100{,}000
(95\,\% CI 0.8--2.2)}, concentrated in respiratory (\textbf{0.9, 0.4--1.4}) and
cardiovascular (\textbf{0.6, 0.1--1.1}) subgroups, and driven by the oldest age strata.
Wind direction satisfies exclusion only if it affects COVID outcomes solely through
PM$_{2.5}$ --- doubtful during a period when weather drove both mobility and indoor
crowding --- and the estimates are ecological rate differences, not individual risks.
Still the strongest identification strategy this literature has produced.}
{10.1093/ije/dyag152}

\paperentry{de Souza et al. 2026}
{Wildfire smoke, PM$_{2.5}$ and respiratory morbidity: spatial mapping and time series}
{Sci Total Environ}
{\textbf{1{,}467{,}172 respiratory hospital admissions} in S\~ao Paulo State, 2018--2025
(DATASUS), against \textbf{20{,}446 INPE-recorded wildfire events} and daily CETESB
PM$_{2.5}$, temperature and humidity, analysed by GAM on weekly aggregates with adjustment
for meteorology, calendar month and a nonlinear trend. Children 0--10 and males carried
the greatest burden; population-standardised rates were highest in the north and
north-west, coinciding with fire hotspots, while absolute counts concentrated in the
metropolitan area --- a textbook illustration of why per-capita and count maps disagree.
\textbf{Each 10\,$\mu$g/m$^3$ weekly PM$_{2.5}$ was associated with a 3.4\,\% increase in
admission rates (RR 1.034, 95\,\% CI 0.889--1.204, $p = 0.054$)}, strongest at zero lag.
Fire counts and PM$_{2.5}$ correlated at $\rho = 0.732$. The GAM explains 82.8\,\% of
deviance, but weekly aggregation of an acute exposure throws away most of the contrast
and is the likeliest reason the interval spans the null.}
{10.1016/j.scitotenv.2026.182227}

\band{Reproductive \& developmental\hfill 1 record}{Sky}

\paperentry{Chen et al. 2026g}
{Prenatal size-fractionated PM and fetal growth trajectories, mediated by maternal glucose}
{Toxics}
{\textbf{616 mother--fetus pairs} in Guangzhou with repeated mid-to-late-pregnancy
ultrasound (HC, AC, FL, BPD as z-scores; EFW trajectories classed average, slow or rapid)
and residential PM$_1$, PM$_{2.5}$, PM$_{10}$ from 1\,km spatiotemporal models. Each
5\,$\mu$g/m$^3$ increase raised the odds of a \emph{rapid} EFW trajectory versus average:
\textbf{PM$_1$ OR 1.56 (1.03--2.37), PM$_{2.5}$ 1.44 (1.10--1.88), PM$_{10}$ 1.28
(1.05--1.56)} --- monotone in the inverse of particle size, as a deposition argument would
predict. BPD, HC and AC z-scores moved in the same direction. \textbf{Maternal fasting
glucose and insulin resistance mediated 5.0--29.4\,\%} of the association. The mediation
range is wide enough to be uninformative at its lower end, mid-pregnancy glucose is a
single measurement, and \emph{accelerated} fetal growth is a less established adverse
outcome than restriction, so the clinical framing is the weakest part. The size gradient
is the citable result.}
{10.3390/toxics14080709}

\band{Mechanistic toxicology\hfill 3 records}{Violet}

\paperentry{Qin et al. 2026b}
{Engineered protein coronas: ApoE enhances particulate transport across the blood-brain barrier}
{Environ Health (Wash)}
{Gold nanoparticles functionalised with thiolated DNA aptamers serve as a
\textbf{compositionally controlled PM surrogate}, allowing reproducible assembly of
coronas from three physiologically relevant proteins --- human serum albumin,
apolipoprotein E and complement C3b --- and isolating corona identity from particle
identity. In a bEnd.3 Transwell BBB model, \textbf{ApoE-enriched coronas nearly doubled
apparent permeability and increased 24 h transport 1.6-fold} over bare particles, through
an energy-dependent, predominantly clathrin-mediated route. ApoE precoronation also
raised endothelial penetration of polystyrene particles and \textbf{increased the
endothelial cytotoxicity of real industrial-source PM samples}, which is the step that
connects the model to ambient aerosol. The surrogate is the limitation: gold
nanoparticles share neither the surface chemistry nor the morphology of combustion PM,
and a Transwell monolayer is not a neurovascular unit. Mechanistically the cleanest
account yet of why particle \emph{surface}, not particle mass, governs neural access.}
{10.1021/envhealth.5c00803}

\paperentry{Moonwiriyakit et al. 2026}
{A STAT6 reporter assay for screening toxicants affecting allergic airway inflammation}
{J Vis Exp}
{A methods paper establishing a stable STAT6-responsive luciferase reporter in
\textbf{16HBE14o- human bronchial epithelial cells} by lentiviral transduction, with
puromycin selection, MOI optimisation and monoclonal isolation by limiting dilution.
Responsiveness was validated against IL-4/IL-13, pathway specificity confirmed with
selective STAT6, STAT5 and STAT3 inhibitors, and reproducibility quantified by
\textbf{Z$'$-factor}. Applied to individual air-pollution constituents,
\textbf{benzo[b]fluoranthene and particulate matter both significantly increased STAT6
reporter activity}. No dose--response, no PM characterisation (source, size, mass
loading) and no Z$'$ value are given in the abstract, so the assay's screening
sensitivity cannot be assessed from it. The value is infrastructural: allergic-airway
signalling is a plausible route for PM-associated asthma incidence and there has been no
scalable prioritisation assay for it.}
{10.3791/71694}

\paperentry{Nikaido-Landry et al. 2026}
{EGFR dysregulation and climate change: pollution-promoted lung cancer}
{Philos Trans R Soc Lond B Biol Sci}
{Discussion-meeting review connecting rising anthropogenic fine PM to lung cancer in
never-smokers via EGFR signalling. The mechanistic core is the promotion model:
\textbf{fine PM promotes malignant transformation in cells that already carry an EGFR
mutation}, so the relevant population is people with oncogenic mutations in
histologically normal lung tissue rather than people accumulating new mutations from
exposure. That reframing matters --- it predicts a susceptible subgroup rather than a
population-wide dose--response, and it implies precision-prevention strategies targeted
by genotype rather than universal exposure reduction. A narrative review with no
systematic search and no new data; the promotion model rests principally on one line of
murine and cohort work, and the ``climate change'' framing in the title is rhetorical
rather than analytic. Tier B for the mechanistic argument.}
{10.1098/rstb.2024.0508}

\band{Occupational \& indoor\hfill 3 records}{Clay}

\paperentry{Li \& Zheng 2026}
{Size-resolved PM$_{10}$-bound PFAS: respiratory deposition and oxidative stress in recycling workers}
{Toxics}
{Panel study at a southern-Chinese waste recycling plant: \textbf{15 size-resolved PM
samples and 280 repeated first-void urines from 20 workers over 45 consecutive days},
with linear mixed models linking size-resolved PM-bound PFAS to urinary PFAS and to
8-OHdG and MDA. Short-chain PFAS predominated in both matrices; urinary ICCs were fair to
excellent ($>$0.4). \textbf{PFAS in PM$_{2.1-10}$ correlated more strongly with internal
dose than PFAS in PM$_{2.1}$} --- the opposite of the fine-fraction-dominates prior, and
consistent with mucociliary clearance and gastrointestinal uptake of coarse-fraction
deposits. \textbf{Each unit increase in ln average daily intake was associated with
10.2--12.6\,\% higher 8-OHdG (PM$_{4.7-5.8}$, PM$_{9.0-10}$) and 8.0--14.6\,\% higher MDA
(PM$_{5.8-9.0}$)}. Twenty workers, one plant, and the authors call the biomarker
associations suggestive; with this many size bins the multiplicity is substantial and no
correction is described. The size-resolved design is nonetheless the right one and is
almost never done.}
{10.3390/toxics14080682}

\paperentry{Sadki et al. 2026}
{Real-time formaldehyde exposure and acute symptoms in healthcare workers}
{Cad Saude Publica}
{Cross-sectional study of \textbf{180 healthcare workers across 12 hospital departments}
in Morocco, with airborne formaldehyde measured over full eight-hour shifts by calibrated
semiconductor dosimeter alongside a structured symptom questionnaire. \textbf{Every mean
concentration was below the Moroccan occupational limit of 0.25\,mg/m$^3$}
(0.048--0.196\,mg/m$^3$, 19--78\,\% of the limit), yet highly exposed workers reported
markedly more eye irritation (72.5\,\%), headache (63\,\%) and respiratory discomfort
(56.5\,\%). \textbf{Prolonged exposure ($\geq$8 h/day, $\geq$3 days/week) independently
predicted multiple symptoms (AOR 3.2, 95\,\% CI 1.7--5.9)}, as did inadequate ventilation
(2.4) and inconsistent PPE use (1.9); years of exposure correlated with symptom count
($\rho = 0.41$). \textbf{The exposure is gaseous, not particulate}, so the AORs are
excluded from this issue's forest plot. Carried because the deployment pattern ---
low-cost real-time dosimetry as routine occupational surveillance in a resource-limited
setting --- is exactly the use case for distributed sensing, and because symptom
self-report against measured exposure is open to reporting bias in the obvious direction.}
{10.1590/0102-311XEN191325}

\paperentry{Perna et al. 2026}
{Indoor air pollution as a determinant of airport occupational exposure}
{Life (Basel)}
{Narrative review arguing that airport occupational exposure assessment has been anchored
on aircraft engine emissions and has under-weighted \textbf{indoor air quality driven by
indoor--outdoor exchange and ventilation}. Ultrafine particles, black carbon, PAHs and
trace metals are identified as the constituents driving oxidative stress, inflammation
and xenobiotic pathway activation, with effects framed beyond the respiratory system. The
operationally useful recommendation is to monitor indoor microclimatic parameters
--- temperature, pressure, humidity --- as \textbf{event detectors for indoor air-quality
deterioration}, which is a cheap addition to any fixed deployment. No systematic search,
no quantitative synthesis, no exposure or effect estimates; the closing suggestion that
AI-based exposure modelling is promising is asserted with the authors' own caveat that it
lacks validation. Tier C: useful as a framing document for airport deployments, not as
evidence.}
{10.3390/life16081242}

\band{Burden, policy \& mitigation\hfill 5 records}{Coral}

\paperentry{\"Ozmen et al. 2026}
{Long-term PM$_{2.5}$ and premature mortality across Turkiye}
{Toxics}
{Province-level cross-section of \textbf{all 81 Turkish provinces} for 2022, combining
annual mean PM$_{2.5}$ with WHO AirQ+ attributable-death estimates for adults 30+ at a
5\,$\mu$g/m$^3$ counterfactual. \textbf{Population-weighted mean PM$_{2.5}$ was
27.0\,$\mu$g/m$^3$ --- 5.4$\times$ the WHO guideline --- and all 81 provinces exceeded
it.} Bivariate OLS explained 41\,\% of between-province variation in attributable
mortality (\textbf{slope 3.23 additional deaths per 100{,}000 per 1\,$\mu$g/m$^3$,
95\,\% CI 2.37--4.10}); a log-log model gave an elasticity of \textbf{0.80
(0.65--0.95)}. The attributable fraction of natural-cause mortality rose from
\textbf{8.8\,\% in the lowest exposure quartile to 24.6\,\% in the highest}, totalling
\textbf{68{,}440 premature deaths, 14.2\,\% of natural-cause mortality}. The authors are
explicit about the circularity and it is worth repeating: attributable deaths are
\emph{generated} by AirQ+ from a fixed concentration--response function, so regressing
them on concentration recovers the assumed function, not independent evidence. The
burden magnitude stands; the gradient is arithmetic.}
{10.3390/toxics14080728}

\paperentry{Sommar et al. 2026}
{Health economics of PM$_{2.5}$ from residential wood combustion, northern Sweden}
{Int J Environ Res Public Health}
{Health impact assessment for \textbf{280{,}742 residents of V\"asterbotten County} using
high-resolution 2023 dispersion-modelled PM$_{2.5}$ from small-scale residential heating,
population-weighted to home-address coordinates, with the VALESOR framework and a
near-source relative risk of \textbf{1.26 per 10\,$\mu$g/m$^3$} for natural-cause
mortality. \textbf{Mean RWC-attributable exposure was only 0.16\,$\mu$g/m$^3$}
(IQR 0.10--0.21), yielding \textbf{30 years of life lost annually} plus small morbidity
counts (0.3 lung cancers, 4 childhood asthma, 3 COPD, 3 IHD, 1 stroke, 2 dementia) and a
total annual cost of \textbf{$\sim$EUR 7.0 million (EUR 4.6--9.4 million)}. Two
assumptions carry the result: the near-source RR of 1.26 is far above the standard
$\sim$1.08 per 10\,$\mu$g/m$^3$ and is a deliberate wood-smoke-toxicity premium that the
evidence base supports only weakly, and the population-mean of 0.16\,$\mu$g/m$^3$ conceals
a highly skewed near-source distribution. Fractional case counts at this precision are
not meaningful.}
{10.3390/ijerph23081053}

\paperentry{Zhai et al. 2026}
{Global disparities in lifetime lung cancer risk: development, smoking and fine PM}
{Front Public Health}
{GLOBOCAN 2022 with the adjusted-for-multiple-primaries method, giving \textbf{global
lifetime risk of lung cancer incidence 4.91\,\% and mortality 4.07\,\%}, ranging from
North America (7.30\,\% / 5.38\,\%) to Africa (0.99\,\% / 0.93\,\%) and higher in males
on every continent. In country-level regression, smoking prevalence predicted both
(LRI $\beta = 0.001$, $p < 0.001$), and \textbf{PM$_1$ was the strongest environmental
predictor (LRI $\beta = 0.40$, $p = 0.001$; LRM $\beta = 0.20$, $p = 0.006$)}, slightly
ahead of PM$_{2.5}$ (0.36 and 0.18). Higher national development correlated with higher
lifetime risk. The ecological confounding is total: HDI, smoking history, registry
completeness and diagnostic capacity all covary with PM at country level, and lung cancer
incidence in high-HDI countries is dominated by historical smoking and by the ability to
detect it. \textbf{PM$_1$ outranking PM$_{2.5}$ is the one result worth carrying}, and
even that is weak given how few countries have independent PM$_1$ products.}
{10.3389/fpubh.2026.1842037}

\paperentry{Acat et al. 2026}
{National trends in ambient pollutants and cause-specific mortality, Turkiye 2018--2024}
{Atmosphere}
{Seven annual national observations for PM$_{10}$, SO$_2$, CO, NO$_2$ and O$_3$, each the
unweighted mean of eligible station annual means from official bulletins, against crude
national mortality rates, tested by OLS with Mann--Kendall, Sen's slope and a
strict common-station sensitivity analysis. \textbf{PM$_{10}$ fell from 50.4 to
41.0\,$\mu$g/m$^3$} and SO$_2$ from 13.2 to 8.5, both significant. The all-station NO$_2$
indicator rose from 29.2 to 34.4 but \textbf{the trend vanished under common-station
restriction} (slope $-0.021$, $p = 0.950$) --- a clean demonstration that a national mean
over a changing monitoring network measures the network as much as the air.
\textbf{PM$_{2.5}$ was excluded outright} for insufficient reporting comparability, which
is itself the story for a country whose modelled PM$_{2.5}$ burden the AirQ+ record above
puts at 68{,}440 deaths. Only neoplasm mortality declined significantly. Crude national
rates against unweighted station means; no causal reading is possible and none is claimed.}
{10.3390/atmos17090832}

\paperentry{Zhu et al. 2026}
{Short-term effectiveness of urban road cleaning for fugitive dust, from TRAKER mobile monitoring}
{Atmos Pollut Res}
{\textbf{Metadata-only deposit.} TRAKER (Testing Re-entrained Aerosol Kinetic Emissions
from Roads) infers road-dust emission potential from the difference between behind-tyre
and background aerosol measured on a moving vehicle, which makes it one of the few
mobile methods that measures an \emph{emission} rather than a concentration. The question
in the title --- how long the benefit of a street-sweeping pass actually persists --- is
one municipalities repeatedly fund cleaning programmes without answering, and the honest
prior from the existing literature is ``hours to a few days, if at all''. No abstract, no
site, no numbers. Carried at tier C and retained because the method class is directly
relevant to mobile PM sensing.}
{10.1016/j.apr.2026.103190}

\band{Other clinical endpoints\hfill 4 records}{Slate}

\paperentry{Huang et al. 2026}
{PM$_{2.5}$-bound metals and female breast cancer incidence, New Mexico}
{Air Qual Atmos Health}
{Census-tract analysis of New Mexico Tumor Registry breast-cancer incidence, 1997--2019,
against annual PM$_{2.5}$-bound metal exposure intensities from a calibrated Emission
Weighted Proximity Model over a \textbf{10-year pre-diagnosis window}, fitted by GEE with
Gaussian estimation. All five metals were associated with higher age-adjusted incidence:
\textbf{chromium RR 1.07 (1.05--1.08), cobalt 1.08 (1.05--1.11), copper 1.06
(1.04--1.07), manganese 1.05 (1.02--1.08), mercury 1.05 (1.03--1.07)}, with a monotone
trend across unexposed/low/medium/high categories for every metal. The intervals are
strikingly narrow, which is the tell: these are tract-level rate ratios on 23 years of
registry data, so precision reflects the number of tract-years rather than the number of
informative exposure contrasts, and the five metals are emitted by shared sources and are
almost certainly collinear --- five nominally independent estimates from one exposure
surface. EWPM is a proximity model, not a measurement. Ecological throughout.}
{10.1007/s11869-026-02079-1}

\paperentry{Zhang et al. 2026f}
{Plasma metabolomic signatures of air pollution and rheumatoid arthritis risk}
{Front Public Health}
{\textbf{403{,}332 RA-free UK Biobank participants} followed a median \textbf{14.01
years}, with residential annual PM$_{2.5}$, PM$_{10}$, NO$_2$, NO$_x$ and a composite air
pollution score, plasma metabolomics, machine-learning-derived pollutant-specific
metabolic signatures, Cox models, generalised propensity score balancing and mediation
analysis. Per SD of signature: \textbf{PM$_{2.5}$-related HR 1.208 (1.180--1.236)},
NO$_2$ 1.141, NO$_x$ 1.174, APS 1.184, stable under propensity balancing.
\textbf{Signatures accounted for 15.7--21.1\,\%} of the pollution--RA associations.
The critical reading: the hazard ratios are for the \emph{signature}, a metabolite score
trained to predict exposure, not for exposure itself --- a signature can outperform its
own target because it absorbs correlates of exposure (diet, adiposity, smoking) that the
pollutant variable does not. Single-timepoint metabolomics, UK Biobank's healthy-volunteer
selection, and a mediation framework that assumes no exposure--mediator confounding.}
{10.3389/fpubh.2026.1901951}

\paperentry{Chen et al. 2026h}
{Metabolic perturbation in elderly adults exposed to PM$_{2.5}$-bound ester compounds}
{Metabolites}
{\textbf{258 Fuzhou residents aged 60+} with \textbf{72-hour personal PM$_{2.5}$ sampling
(UPAS V2)}, GC-MS speciation and LC-MS plasma metabolomics, targeting non-phthalate,
non-organophosphate esters that routine speciation ignores. Mean personal PM$_{2.5}$ was
\textbf{38.06\,$\mu$g/m$^3$} with esters dominating the measured organic composition;
20 high-concentration esters were prioritised. Covariate-adjusted linear regression
associated ester exposure with \textbf{steroid biosynthesis, glycolysis/gluconeogenesis,
glycerophospholipid and purine/pyrimidine} pathway alterations. Personal rather than
ambient sampling in an elderly cohort is the strength and is uncommon; against that, the
ester annotations are putative, the design is cross-sectional with a single 72-h exposure
window against a single metabolome, pathway enrichment on untargeted features has a
well-known false-discovery problem, and the authors state the work is hypothesis-
generating. No effect sizes are reportable from the abstract.}
{10.3390/metabo16080553}

\paperentry{Zinan'dala et al. 2026}
{Urinary metabolomic alterations across contrasting ambient pollution regions, Thailand}
{Int J Mol Sci}
{Untargeted $^1$H-NMR urinary metabolomics comparing \textbf{51 adults in Chiang Mai
(high exposure) with 52 in Songkhla (low exposure)} during a regional pollution episode,
using long-term residence as the exposure proxy. PLS-DA separated the groups
(\textbf{$R^2 = 0.873$ but $Q^2 = 0.447$} --- good fit, modest prediction, and the gap is
the standard warning sign of overfitting in this method). Eight metabolites differed at
$p < 0.05$, six surviving Benjamini--Hochberg; \textbf{L-tryptophan, hippuric acid,
xanthine and 5-HIAA remained significant after covariate adjustment}, with L-arginine and
L-cystathionine lower in the high-exposure group. The design's fatal ambiguity is
acknowledged by the authors: \textbf{region is the exposure variable}, and Chiang Mai and
Songkhla differ in diet, water, altitude and season as much as in PM$_{2.5}$. Hippuric
acid in particular is a dietary polyphenol marker before it is anything else. Tier B for
method, not for attribution.}
{10.3390/ijms27167193}

\clearpage

% ==================================================================== TABLE
\section*{\sffamily\bfseries\color{Deep}Corpus register}
\vspace{-1.5mm}{\color{Amber}\rule{\linewidth}{1.1pt}}\vspace{2.5mm}

{\fontsize{6.7}{8.4}\selectfont
\setlength{\tabcolsep}{2.4pt}
\begin{longtable}{@{}>{\RaggedRight}p{25mm}>{\RaggedRight}p{41mm}>{\RaggedRight}p{27mm}%
>{\RaggedRight}p{23mm}>{\RaggedRight}p{18mm}>{\RaggedRight}p{15mm}@{}}
\toprule
\textsf{\textbf{First author}} & \textsf{\textbf{Focus}} & \textsf{\textbf{Journal}} &
\textsf{\textbf{Design}} & \textsf{\textbf{Metric}} & \textsf{\textbf{Region}}\\
\midrule
\endfirsthead
\toprule
\textsf{\textbf{First author}} & \textsf{\textbf{Focus}} & \textsf{\textbf{Journal}} &
\textsf{\textbf{Design}} & \textsf{\textbf{Metric}} & \textsf{\textbf{Region}}\\
\midrule
\endhead
\multicolumn{6}{@{}l@{}}{\cellcolor{Mist}\textsf{\textbf{\textcolor{Deep}{Neuro / mental health}}}}\\[0.6mm]
\href{https://doi.org/10.1001/jamanetworkopen.2026.25015}{Dales} & Wildfire-sourced PM$_{2.5}$ and ED visits for migraine and other primary headache & \textit{JAMA Netw Open} & Case-crossover & PM$_{2.5}$ (wildfire) & Canada\\
\href{https://doi.org/10.1093/ageing/afag215}{Glover} & Umbrella meta-review of urban environmental exposures and cognitive health & \textit{Age Ageing} & Umbrella review / meta-analysis & PM$_{2.5}$ & Global\\
\href{https://doi.org/10.1016/j.nsa.2026.107019}{Goossens} & In utero air pollution exposure and pre-/postnatal brain development & \textit{Neurosci Appl} & Narrative review & PM$_{2.5}$ / PM$_{10}$ / PM$_{\mathrm{coarse}}$ & Global\\
\href{https://doi.org/10.1016/j.isci.2026.116617}{Ma} & PM$_{2.5}$-induced neurotoxicity and depression: exposure risks and mechanisms & \textit{iScience} & Narrative review & PM$_{2.5}$ + components & Global\\
\href{https://doi.org/10.3760/cma.j.cn112150-20250804-00752}{Pan} & Maternal PM exposure and offspring neurodevelopment: effects and mechanisms & \textit{Zhonghua Yu Fang Yi Xue Za Zhi} & Narrative review & PM (bulk) & China\\
\addlinespace[0.9mm]
\multicolumn{6}{@{}l@{}}{\cellcolor{Mist}\textsf{\textbf{\textcolor{Teal}{Cardiovascular \& metabolic}}}}\\[0.6mm]
\href{https://doi.org/10.3760/cma.j.cn112338-20251120-00843}{Huang} & PM$_{2.5}$ and dyslipidaemia in rural Chinese children; dietary effect modification & \textit{Zhonghua Liu Xing Bing Xue Za Zhi} & Cross-sectional (multistage) & PM$_{2.5}$ & China\\
\href{https://doi.org/10.32394/pe/218966}{Kazmierska} & Air pollution and cardiovascular disease: review of current data & \textit{Przegl Epidemiol} & Narrative review & PM$_{2.5}$ / PM$_{10}$ & Global\\
\href{https://doi.org/10.1007/s00125-026-06814-2}{Xi} & Air pollution, built environment and temperature vs. prediabetes phenotypes & \textit{Diabetologia} & Prospective cohort & PM$_{2.5}$ & Germany\\
\href{https://doi.org/10.1161/JAHA.125.046917}{Yang} & Short-term PM$_{2.5}$ and hospitalisation/ED use in atrial fibrillation & \textit{J Am Heart Assoc} & Case-crossover & PM$_{2.5}$ (wildfire-influenced) & USA\\
\addlinespace[0.9mm]
\multicolumn{6}{@{}l@{}}{\cellcolor{Mist}\textsf{\textbf{\textcolor{Sky}{Reproductive \& developmental}}}}\\[0.6mm]
\href{https://doi.org/10.1016/j.jes.2025.10.043}{Chen} & Prenatal air pollution, fetal growth trajectory phenotypes and child neurodevelopment & \textit{J Environ Sci (China)} & Prospective cohort & PM$_{2.5}$ / PM$_{10}$ & China\\
\href{https://doi.org/10.1016/j.ecoenv.2026.120566}{Huang} & Joint prenatal bisphenol + air-pollutant mixture exposure and preschool IQ & \textit{Ecotoxicol Environ Saf} & Prospective cohort (mixtures) & PM$_{2.5}$ / PM$_{10}$ & China\\
\href{https://doi.org/10.1007/s40572-026-00552-8}{Peterson} & Ambient air pollution, PAHs and pubertal development: critical review & \textit{Curr Environ Health Rep} & Critical review & PM$_{2.5}$ components & Global\\
\href{https://doi.org/10.1016/j.envres.2026.125313}{Xie} & Long-term PM$_{2.5}$ and timing of pubertal development in boys and girls & \textit{Environ Res} & Prospective cohort & PM$_{2.5}$ & China\\
\addlinespace[0.9mm]
\multicolumn{6}{@{}l@{}}{\cellcolor{Mist}\textsf{\textbf{\textcolor{Sage}{Respiratory \& allergic}}}}\\[0.6mm]
\href{https://doi.org/10.1007/s00484-026-03289-z}{Bayar Muluk} & Climate change and otolaryngology: airway disease impacts and emission reduction & \textit{Int J Biometeorol} & Narrative review & PM$_{2.5}$ & Global\\
\href{https://doi.org/10.4103/mgr.MEDGASRES-D-25-00164}{Liao} & Short-term air pollution, respiratory symptoms and lung function in asthma & \textit{Med Gas Res} & Panel study & PM$_{2.5}$ / PM$_{10}$ & China\\
\href{https://doi.org/10.3390/epidemiologia7040097}{Sanchez de Prada} & Ambient PM and asthma-coded admissions in children 0-3 y, Spain & \textit{Epidemiologia} & Ecological panel & PM$_{2.5}$ / PM$_{10}$ & Spain\\
\href{https://doi.org/10.1016/j.jhazmat.2026.143076}{Yang} & Allergenic vs. toxic risk of airborne seed hairs compared with ambient PM$_{2.5}$ & \textit{J Hazard Mater} & Field measurement & PM$_{2.5}$ + bioaerosol & China\\
\addlinespace[0.9mm]
\multicolumn{6}{@{}l@{}}{\cellcolor{Mist}\textsf{\textbf{\textcolor{Violet}{Mechanistic toxicology}}}}\\[0.6mm]
\href{https://doi.org/10.1016/j.jes.2025.12.074}{Ding} & PM$_{2.5}$ promotes vascular calcification via miR-27b-5p/HES1 signalling & \textit{J Environ Sci (China)} & In vitro + organoid + murine & PM$_{2.5}$ & China\\
\href{https://doi.org/10.1007/s00360-026-01698-5}{Maraschi} & Osmoregulatory disruption and metal bioaccumulation from settleable atmospheric PM & \textit{J Comp Physiol B} & Controlled exposure (ecotox) & Settleable APM & Brazil\\
\href{https://doi.org/10.1016/j.jes.2025.11.018}{Wu} & Macrophage NLRP3 inflammasome in real-ambient PM-induced hepatic glucose dysregulation & \textit{J Environ Sci (China)} & Real-ambient murine + co-culture & Real-ambient PM & China\\
\addlinespace[0.9mm]
\multicolumn{6}{@{}l@{}}{\cellcolor{Mist}\textsf{\textbf{\textcolor{Clay}{Occupational \& indoor}}}}\\[0.6mm]
\href{https://doi.org/10.3389/fmedt.2026.1796828}{Chen} & Nested suction instrument to cut intraoperative surgical-smoke PM$_{2.5}$ & \textit{Front Med Technol} & Bench / device experiment & PM$_{2.5}$ (surgical smoke) & China\\
\href{https://doi.org/10.1093/annweh/wxag055}{Cole} & Determinants of respirable crystalline silica exposure in Australian tunnelling & \textit{Ann Work Expo Health} & Observational exposure study & Respirable dust / RCS & Australia\\
\href{https://doi.org/10.1016/j.jhazmat.2026.143081}{Ding} & Stage-resolved hair Zn-isotope dynamics as a firefighting combustion-exposure tracer & \textit{J Hazard Mater} & Repeated-measures biomarker & Combustion-derived PM & China\\
\href{https://doi.org/10.3390/jox16040119}{Ramsden} & Xenobiotic hazards in aircraft cabin air & \textit{J Xenobiot} & Review + risk economics & Ultrafine particles & Global\\
\addlinespace[0.9mm]
\multicolumn{6}{@{}l@{}}{\cellcolor{Mist}\textsf{\textbf{\textcolor{Amber}{Sensing, forecasting \& instrumentation}}}}\\[0.6mm]
\href{https://doi.org/10.1016/j.jenvman.2026.130515}{Shi} & Adaptive spatiotemporal graph transformer for multi-site multi-step PM$_{2.5}$ forecasting & \textit{J Environ Manage} & Deep-learning model & PM$_{2.5}$ & China / India\\
\href{https://doi.org/10.3390/bios16070380}{Tsuchiya} & Seasonal and indoor-outdoor characteristics of fluorescent aerosol particles in offices & \textit{Biosensors} & Real-time sensor campaign & Bioaerosol / coarse & Japan\\
\addlinespace[0.9mm]
\multicolumn{6}{@{}l@{}}{\cellcolor{Mist}\textsf{\textbf{\textcolor{Amber}{Exposure assessment \& modelling}}}}\\[0.6mm]
\href{https://doi.org/10.1029/2025GH001621}{Solorzano-Araque} & Multiscale exposure modelling: LOTOS-EUROS + mobility surveys + particle toxicity & \textit{GeoHealth} & CTM + mobility integration & PM$_{2.5}$ / PM$_{10}$ & Colombia\\
\addlinespace[0.9mm]
\multicolumn{6}{@{}l@{}}{\cellcolor{Mist}\textsf{\textbf{\textcolor{Coral}{Burden, policy \& mitigation}}}}\\[0.6mm]
\href{https://doi.org/10.1002/hsr2.72886}{Karamzad} & Ischaemic heart disease burden in MENA 1990-2021; PM attributable fraction 28.2\% & \textit{Health Sci Rep} & GBD secondary analysis & Ambient PM & MENA\\
\href{https://doi.org/10.1371/journal.pone.0347989}{Rees} & HEPA filtration to reduce respiratory infection in care homes & \textit{PLoS One} & RCT process evaluation & Filtered PM & UK\\
\href{https://doi.org/10.1016/j.jes.2026.01.009}{Shan} & Policy-driven emission reductions and post-pandemic rebound, high-resolution inventories & \textit{J Environ Sci (China)} & Emission inventory & PM emissions & China\\
\addlinespace[0.9mm]
\multicolumn{6}{@{}l@{}}{\cellcolor{Mist}\textsf{\textbf{\textcolor{Slate}{Other clinical endpoints}}}}\\[0.6mm]
\href{https://doi.org/10.3390/dj14070440}{Lupita} & PM$_{2.5}$, preventive-care capacity and caries/edentulism burden across the EU & \textit{Dent J} & Ecological panel & PM$_{2.5}$ & EU-27\\
\href{https://doi.org/10.1016/j.jdin.2026.06.008}{Park} & Long-term PM exposure and risk of androgenetic alopecia & \textit{JAAD Int} & Nationwide cohort & PM$_{2.5}$ / PM$_{10}$ & South Korea\\
\href{https://doi.org/10.3171/2026.3.PEDS2615}{Sun} & Environmental exposure and area deprivation vs. paediatric hydrocephalus outcomes & \textit{J Neurosurg Pediatr} & Registry cohort + geospatial & PM$_{2.5}$ (satellite-derived) & USA\\
\addlinespace[0.9mm]
\bottomrule
\end{longtable}}


\vspace{4mm}
\begin{keybox}
{\sffamily\bfseries\footnotesize\color{Deep}Provenance and method}\\[1.2mm]
{\fontsize{7.2}{8.8}\selectfont
Entry window \textbf{2026-08-27 to 2026-08-27}, a single day of deposits, built on the
28 August run under the 22:00 rule.
\textbf{Sources.} PubMed E-utilities on \texttt{[Date - Entry]}, run as two separate
queries (health axis 19, instrumentation axis 5, 2 overlapping). Europe PMC
\texttt{CREATION\_DATE} 31, including one Preprints.org record. Crossref by-ISSN across
the tracked non-PubMed-indexed titles, 43 deposits. \textbf{OpenAlex returned HTTP 429
on top of its existing 403 date-filter failure} and contributed nothing; arXiv
contributed nothing new. A \textbf{Consensus} sweep on the forward-scattering /
refractive-index axis returned only the 2023 antecedent work and no in-window record, so
no backfill section appears in this issue.
\textbf{Screening.} 98 raw records, 95 unique after PMID/DOI/title collapse, 6 already in
\texttt{seen.json}, \textbf{36 retained and 53 rejected} with reason logged. Rejections
were dominated by a single-day \textit{Meas Sci Technol} volume release (13 records, none
about aerosol --- bearing diagnostics, optical flow, ultrasonic weld simulation), four
Elsevier retraction notices, and non-airborne environmental chemistry (soil, sediment,
water, wildlife mercury). Two records were rejected on scope after reading:
\textit{Gaseous air pollutant exposure and atherosclerosis} (IJMS), which carries no
particulate exposure metric, and \textit{Air pollution-induced oxidative stress after
CABG} (Antioxidants), a conceptual perspective with no exposure measurement.
\textbf{Four abstract-free deposits} are carried at tier C on title, journal and method
class alone, and are labelled as such in the digest.
\textbf{Effect estimates.} Thirteen from five studies. Excluded by design: the Guangzhou
toddler-sleep betas (negative hour differences, not drawable on a log-ratio axis), the
Turkiye OLS slope (an absolute rate difference), the Moroccan formaldehyde AORs (a
gas-phase contrast) and the NO$_2$/NO$_x$/APS metabolic-signature hazard ratios (not
particulate contrasts).
\textbf{Harvester change made this run.} Morocco, Pakistan, Indonesia, Reunion, Belgium
and the Arctic were added to the geography substring map in \texttt{plots.py}; all six
were falling through the unmapped diagnostic into \textit{Global / multi-region}.
\textbf{Labels.} Subtopic, design, particle-metric and geography are single-label by
dominant contribution --- a judgement, not a controlled vocabulary. All DOI links were
resolved before build.\par}
\end{keybox}

\end{document}
